% ========================================================
% CONCLUSIONES
% ========================================================
\section{Conclusiones}

\subsection{Conclusiones del grupo}

La pregunta guía de la revisión bibliográfica de la PC1 era cómo afectan la granularidad del trabajo
y la estrategia de reducción al rendimiento y a la equivalencia numérica de K-means sobre más de un
millón de registros. Este entregable la responde con datos.

\textbf{La granularidad importa, pero al revés de lo que se intuye.} Trocear fino no cuesta: mandar
un índice por un canal es despreciable frente a procesar mil viajes, y el óptimo es una meseta ancha.
Lo que castiga es trocear grueso, porque deja workers sin trabajo. La regla es muchos más bloques que
workers, y con ella el reparto por canal equilibra la carga solo.

\textbf{La estrategia de reducción decide la reproducibilidad.} Acumular por bloque y reducir en
orden fijo cuesta unos 66~KB y a cambio hace que el resultado concurrente sea el mismo para
cualquier cantidad de workers, en cualquier máquina. Sin eso, la comparación con la versión
secuencial sería una inspección a ojo. La equivalencia con la secuencial se comprueba con tolerancia,
como anticipó la PC1, y la diferencia medida es del orden del redondeo.

\textbf{El límite no es una sección secuencial, es coordinar.} La fracción serial efectiva crece con
los workers, así que el techo de Amdahl no sirve como predicción; la escalabilidad débil, en cambio,
se sostiene hasta cuatro workers y se degrada donde la coordinación empieza a pesar. Sobre el
dataset del trabajo, \cifra{speedup-p4}$\times$ con cuatro workers y \cifra{speedup-p8}$\times$ con
ocho justifican la concurrencia; por debajo de unos \cifra{n-equilibrio}~viajes no la justificarían.

\textbf{Verificar y probar son cosas distintas, y hacen falta las dos.} Spin recorrió todos los
entrelazados de un modelo reducido y confirmó que ningún punto se pierde ni se duplica y que los
centroides no se escriben mientras se leen; el mutante demostró que el modelo detecta la carrera que
dice detectar. \texttt{go test -race} comprobó lo mismo sobre la implementación, en los entrelazados
que le tocó observar. Ninguna de las dos reemplaza a la otra.

\textbf{Medir bien es una decisión de diseño.} El orden barajado, la media recortada, el cociente de
medias y el bootstrap no son adornos: sin ellos el speedup depende de qué corrió primero y de cuántos
picos del sistema operativo cayeron en cada muestra. El protocolo está en código y los tiempos crudos
se guardan, así que cualquier cifra de este informe se puede recalcular sin volver a medir.

\subsection{Trabajo pendiente para el Entregable 3}

Queda elegir $k$ con un criterio de validación interna sobre una muestra, medir la variante con
\texttt{sync.Mutex} que este entregable retiró, verificar propiedades de progreso en el modelo con
hipótesis de \emph{fairness} declaradas, extender la escala a varios meses y contrastar con una
implementación en GPU, sabiendo que ese contraste compara configuraciones completas y no una sola
variable. Y, sobre todo, interpretar los clusters: hasta acá el trabajo demuestra que el algoritmo
concurrente calcula lo mismo que el secuencial, más rápido; qué tipos de viaje encontró y cómo se
distribuyen por zona es la pregunta del caso, y sigue abierta.
